
\documentclass{article}

\usepackage[utf8]{inputenc}
\usepackage{array}
\usepackage{multirow}
\renewcommand{\arraystretch}{1.2}

\usepackage{amsmath, amssymb, amsfonts}
\usepackage{float}
\usepackage{graphicx}
\usepackage{tikz}
\usepackage{tikz-cd}
\usepackage[labelfont=bf]{caption}
\usepackage{subcaption}
\usepackage{enumitem}
\usepackage{hyperref}
\usepackage{pdflscape}
\usepackage{placeins}
\usepackage{footnote}
\usepackage{url}
\usepackage{natbib}
\makesavenoteenv{table}
\makesavenoteenv{tabular}


\def\Pr{{\mathrm{P}}}
\def\Pdo{{\mathrm{P}^*}}

\DeclareMathOperator{\Anobs}{Anobs}
\DeclareMathOperator{\Paobs}{Paobs}
\DeclareMathOperator{\Chobs}{Chobs}
\DeclareMathOperator{\Disobs}{Disobs}



\tikzset{rv/.style={circle, draw, inner sep=0.5mm, minimum size=7mm}}
\tikzset{fv/.style={draw}}

\newcommand{\bW}{\boldsymbol{W}}
\newcommand{\bU}{\boldsymbol{U}}
\newcommand{\bV}{\boldsymbol{V}}
\newcommand{\bR}{\boldsymbol{R}}
\newcommand{\bS}{\boldsymbol{S}}
\DeclareMathOperator{\Do}{do}
\DeclareMathOperator{\Pa}{Pa}

%\usepackage{lineno}
%\linenumbers


\usepackage{tikz}
\usetikzlibrary{calc,trees,positioning,arrows,chains,shapes.geometric,%
    decorations.pathreplacing,decorations.pathmorphing,shapes,%
    matrix,shapes.symbols}
\tikzstyle{def} = [rectangle, rounded corners, minimum width=1cm, minimum height=1cm,text centered, draw=black]
\tikzstyle{opt} = [rectangle, rounded corners, minimum width=1cm, minimum height=1cm,text centered, text width=3cm, draw=black, dashed]
\tikzstyle{arrow} = [thick,->,>=stealth]
\usepackage{pgfplots}
\pgfplotsset{width=10cm,compat=1.9}
\newcommand\independent{\protect\mathpalette{\protect\independenT}{\perp}}
\def\independenT#1#2{\mathrel{\rlap{$#1#2$}\mkern2mu{#1#2}}}
\newcommand\indep{\independent}


\setlength{\oddsidemargin}{+0 in} \setlength{\evensidemargin}{-0 in}
\setlength{\textwidth}{6.5 in} \setlength{\topmargin}{-.125 in} \setlength{\headsep}{0in}
\setlength{\textheight}{9 in}



\DeclareCaptionFormat{mysize}{\fontsize{8}{9}\selectfont#1#2#3}

\usepackage{mdframed}

\usepackage[flushleft]{threeparttable}
\newcommand{\pkg}[1]{{\fontseries{b}\selectfont #1}}

\def\spacingset#1{\renewcommand{\baselinestretch}%
{#1}\small\normalsize} \spacingset{1.5}


\begin{document}

\begin{center}
{\Large Response to review}\\[0.5em]
``Deriving Complete Constraints in Hidden Variable Models''
\end{center}

\vspace{-0.5cm}
{\sl We would like to thank the associate editor and the referees for their useful comments. We have incorporated their suggestions into the manuscript, and we find it much improved. Please see our point-to-point responses below within the

\begin{mdframed} boxes \end{mdframed}}

{\sl Passages set off as quotations are taken from the revised manuscript and its Appendix. The manuscript has been reorganized, so we use the new numbering throughout. Propositions 1, 2, 3 and 4 of the submitted version are now Propositions 3, 4, 5 and 6. The HLP and Strong Face Splitting propositions, previously Propositions 6 and 7, are now Propositions 1 and 2. Proposition 5 of the submitted version, the rule of \citet{evans2016graphs}, is no longer stated separately. Remark 3 is new, and the remark on merging unobserved variables, previously Remark 2, is now Remark 4. Figures 5 and 7 keep their numbers.}

\section*{Reviewer 1}

\begin{enumerate}
    \item The formula in Proposition 1 can give the wrong interventional distribution by omitting a necessary conditioning variable. Consider $Z \to X \leftarrow H \to Y \leftarrow W$ \ldots. Please correct the conditioning formula and revisit Propositions 2--3, also specifying how conditioning events of probability zero are handled.
    \begin{mdframed}
        The reviewer is correct, and we are grateful for the clear counterexample (which is the bipartite Bell scenario). The error was in the written statement only, which we regret.
        In the submitted version, the accompanying software (\texttt{dhvm\_district\_params}) conditioned each $V_{1i}$ on all variables in the district and its external parents that precede $V_{1i}$ in a topological order of the whole graph. That set is also valid, so the numerical results in the paper are unaffected. We checked each example. The stated formula and the correct one coincide for the IV, sequential IV, front door, Figure 5 and Figure 7 examples. They differ only for the tripartite Bell example, and there the paper already wrote the constraints in terms of $\Pr(v_1 v_2 v_3 \mid xyz)$, which is the correct district factor.

        Proposition 3 now reads:
        \begin{quote}
        Order $\bV_1=(V_{11},\ldots,V_{1k})$ consistently with a topological ordering of $\mathcal{G}$, let $\bV_1^{<i}=(V_{11},\ldots,V_{1,i-1})$ and $\bV_1^{\leq i} = (V_{11},\ldots,V_{1i})$, and let
        \[
        \boldsymbol{C}_{i} = \bV_1^{<i} \cup \left\{\bV_2 \cap \Paobs(\bV_1^{\leq i})\right\},
        \]
        i.e., the earlier members of the district together with the external observed parents of $V_{1i}$ and of every earlier member. Writing $c_i$ for the corresponding subvector of $(v_1, v_2)$, define
        \[
            \Pdo(\bV_1=v_1\mid \bV_2=v_2)
            =
            \prod_{i=1}^k
            \Pr\!\left(
                V_{1i}=v_{1i}
                \,\middle|\,
                \boldsymbol{C}_{i} = c_i
            \right).
        \]
        \end{quote}
        For the reviewer's graph this gives $\Pr(x \mid z)\Pr(y \mid x, z, w)$, the target that the reviewer states. The text after the proposition adds:
        \begin{quote}
        Under the conditional independencies implied by $\mathcal{G}$, conditioning additionally on any other variables that precede $V_{1i}$ in a topological ordering does not change the factor (see the proof of Proposition 3), but $\boldsymbol{C}_i$ cannot be made smaller in general.
        \end{quote}

        The proof of Proposition 3 (Appendix C) now derives $\Pdo$ from Lemma 1 of \citet{tian2002identification}, which conditions each $V_{1i}$ on all variables that precede it, and then reduces these conditioning sets to $\boldsymbol{C}_i$. It uses c-degree one at this step:
        \begin{quote}
        Because $D_1$ has c-degree 1, all members of $D_1$ share the single latent parent, so $T_i = \bV_1^{\leq i}$.
        \end{quote}
        and it addresses the error directly:
        \begin{quote}
        In general, however, $\boldsymbol{C}_i$ cannot be replaced by $\bV_1^{<i} \cup \{\bV_2 \cap \Anobs(V_{1i})\}$. External parents of earlier members of the district become informative about the shared latent variable once those members are conditioned on.
        \end{quote}
        Proposition 4 (causal interpretation) follows from this, and it now has its own short proof in Appendix C.

The proof of Proposition 5 (completeness) previously appealed to the district factorization informally. It now states the key step explicitly:
        \begin{quote}
        It remains to show that $\Pr(\bV = v) = \prod_{j} \Pdo(\bV_1^j = v^j_1 \mid \bV_2^j = v^j_2)$ for every $v$.
        \end{quote}
        It obtains this from the chain rule and the ordered local Markov property, and it concludes:
        \begin{quote}
        Note that this step uses the conditional independencies essentially. Without them, the observed distribution need not equal the product of its district factors.
        \end{quote}

Regarding zero-probability events, the main development now assumes positivity:
        \begin{quote}
        Throughout, we assume that the observed distribution is strictly positive, i.e., $\Pr(\bV = v) > 0$ for all $v$, so that all conditional probabilities below are well defined. The boundary case is discussed in Remark 2.
        \end{quote}
        Remark 2 specifies how conditioning events of probability zero are handled, and it is proved in Appendix C:
        \begin{quote}
        The first factor that is not positive decides the entry of $\boldsymbol{p}$. If it is zero, the entry is zero, whatever follows it. If it is undefined, because $\Pr(\boldsymbol{C}_i = c_i) = 0$, the observed distribution does not determine the entry, and we treat the entry as a free variable. Then $\Pr \in \mathcal{M}(\mathcal{G})$ if and only if $\Pr$ satisfies the conditional independencies and, for each district, the free entries can be chosen so that $H_i \boldsymbol{p} \leq b_i$ and $H_e \boldsymbol{p} = b_e$. Equivalently, there is a probability vector $\boldsymbol{r}$ with $\boldsymbol{p} = \boldsymbol{B}\boldsymbol{r}$ on the remaining entries, which is a linear program.
        \end{quote}
    \end{mdframed}

    \item Listing all observable constraints is a different task from obtaining causal bounds, as in \citet{duarte2023automated}. Which part of the completeness argument requires a new result beyond Bonet's IV construction and existing district-factorization results? \ldots \,
    \begin{mdframed}
    We agree, and we now say so in a new subsection, ``Relation to previous work'' (Section 2.7):
    \begin{quote}
    Each is known on its own. What is new is the following. First, we identify c-degree one as the class of graphs where every district factor is a linear image of a simplex, so the polytope argument applies one district at a time (Proposition 3). Second, we show that the district polytopes together with the ordinary conditional independencies are \emph{sufficient}, i.e., that nothing is lost when the joint distribution is reassembled from its district factors (Proposition 5). This step depends on the conditioning sets in Proposition 3 being chosen correctly and on the ordered local Markov property. Third, we give an exact criterion for which constraints are nontrivial (Proposition 6), and we use it to link the equality constraints to the nested Markov model. We also give a pre-processing step for graphs with higher c-degree and an open-source implementation.
    \end{quote}
    As the reviewer's first comment shows, the second step is sensitive to exactly how $\Pdo$ is defined. Bonet's argument covers only the IV model, which has a single district with one unobserved variable. The same subsection contrasts our aim with that of \citet{duarte2023automated}:
    \begin{quote}
    Our goal also differs from the related work of \citet{duarte2023automated}, who use the same response function polynomials to bound causal effects by optimization. We instead aim to list every restriction that the model places on the observed distribution, without reference to any target quantity.
    \end{quote}
    \end{mdframed}
\item Could the authors instead give a nontrivial inequality among the second district's 1,118 potentially nontrivial constraints, together with a probability table that satisfies the equality constraints but violates that inequality?
\begin{mdframed}
    The reviewer is right that the displayed inequality for the first district is trivial. In fact, the first district imposes no restrictions at all on the observed distribution, and we have rewritten that paragraph accordingly (Section 5):
    \begin{quote}
    For the first district in the graph in Figure 7 that contains $V_1, V_3, V_6$, Proposition 6 finds no nontrivial constraints. Besides normalization, the equality constraints are
    \[
    \Pdo(V_1 = v_1, V_6 = v_6 \mid V_2 = 0) = \Pdo(V_1 = v_1, V_6 = v_6 \mid V_2 = 1), \quad v_1, v_6 \in \{0, 1\},
    \]
    which are of type (O) with $\bS = \{V_1, V_6\}$. The district factor is $\Pdo(v_1, v_3, v_6 \mid v_2) = \Pr(v_1)\Pr(v_6 \mid v_1)\Pr(v_3 \mid v_1, v_6, v_2)$, so $\Pdo(v_1, v_6 \mid v_2) = \Pr(v_1, v_6)$. The equalities therefore hold identically, and the inequalities reduce to nonnegativity.
    \end{quote}

For the second district, we now give an explicit nontrivial example as the referee suggested. The figure of 1,118 was the number of constraints flagged by the screening procedure of the submitted version, which did not use the equality constraints. With the exact criterion of Proposition 6 (see comment 5 below), the count changes:
    \begin{quote}
    After merging $U_1$ and $U_3$ in the second district and applying our algorithm, Proposition 6 identifies 1096 nontrivial inequality constraints and 6 linearly independent nontrivial equality constraints.
    \end{quote}
    The nontrivial equality constraints are described as follows:
    \begin{quote}
    In observed terms this reads
    \[
    \sum_{v_2} \Pr(v_2 \mid v_1)\Pr(v_4, v_5 \mid v_1, v_2, v_3) \mbox{ does not depend on } v_1,
    \]
    which is a Verma-type constraint, and which a randomly chosen distribution violates. Together with the trivial equality constraints, this condition implies all 6 nontrivial equality constraints of the district.
    \end{quote}
    The manuscript then gives an inequality implied by the constraints, and a distribution that satisfies the equality constraints but violates it:
    \begin{quote}
    The inequality constraints imply
    \begin{equation}
    \sum_{v_1, v_3 \in \{0,1\}} \; \sum_{v_2, v_4, v_5 \in \{0,1\}} a(v_1, v_2, v_3, v_4) \, \Pdo(v_2, v_4, v_5 \mid v_1, v_3) \leq 3, \tag{3}
    \end{equation}
    where $a(v_1, v_2, v_3, v_4) = (1 - v_1 v_3)\,\mathbf{1}\{v_2 = v_4\} + v_1 v_3\,\mathbf{1}\{v_2 \neq v_4\}$ and $\mathbf{1}\{\cdot\}$ is the indicator function. Thus $a = 1$ if $V_2$ and $V_4$ are equal when $v_1 v_3 = 0$, or differ when $v_1 = v_3 = 1$, and $a = 0$ otherwise. For each $(v_1, v_3)$, the inner sum in (3) is the probability of this event under $\Do(V_1 = v_1, V_3 = v_3)$. The inequality states that these four probabilities sum to at most 3. This is a Bell-type inequality. To see that it is nontrivial, consider the following distribution of binary variables:
    \begin{itemize}
        \item $\Pr(V_1 = 1) = 1/2$, and $\Pr(V_2 = 1 \mid v_1) = \Pr(V_6 = 1 \mid v_1) = \Pr(V_3 = 1 \mid v_1, v_6, v_2) = \Pr(V_5 = 1 \mid v_1, v_2, v_3, v_4) = 1/2$ for all values of the conditioning variables;
        \item $\Pr(V_4 = v_4 \mid v_1, v_2, v_3) = 1/8 + \frac{3}{4}\, a(v_1, v_2, v_3, v_4)$, that is, $7/8$ if $a(v_1, v_2, v_3, v_4) = 1$ and $1/8$ otherwise.
    \end{itemize}
    Every probability in the resulting joint distribution of $(V_1, \ldots, V_6)$ is at least $1/256$. The distribution satisfies all the conditional independencies and all the equality constraints of both districts; in particular, $\Pdo(V_4 = 1 \mid v_1, v_3) = 1/2$ for all $v_1, v_3$. Yet for each of the four values of $(v_1, v_3)$, the inner sum in (3) equals $7/8$, so the left side equals $7/2 > 3$. By Remark 4, this distribution is therefore also incompatible with the original graph in Figure 7, left.
    \end{quote}
    The manuscript also gives the shortest of the returned constraints that this distribution violates:
    \begin{quote}
    The example distribution violates 13 of the inequality constraints returned for this district. The shortest of them is
    \begin{align*}
    \Pdo(V_2 = 0, V_4 = 1, V_5 = 0 \mid V_1 = 1, V_3 = 1) \leq{} & \Pdo(V_2 = 1, V_4 = 0 \mid V_1 = 0, V_3 = 0) \\
    & + \Pdo(V_2 = 0, V_4 = 1, V_5 = 0 \mid V_1 = 1, V_3 = 0) \\
    & + \Pdo(V_2 = 0, V_4 = 1, V_5 = 0 \mid V_1 = 0, V_3 = 1).
    \end{align*}
    For the example distribution, the left side equals $7/32$ and the right side equals $1/16 + 1/32 + 1/32 = 1/8$, so the constraint is violated.
    \end{quote}
\end{mdframed}

    \item \ldots\, Please report matrix dimensions, constraint counts and runtimes for several such cases.
    \begin{mdframed}
        We have added Table 1 in Section 6. It covers eight districts, including instruments with three and four levels and the tripartite Bell district. The times are for a single laptop, as stated in the caption, and the benchmarks are reproduced by the script \texttt{10-benchmarks.R} in the supplementary material. The text explains how the cost grows:
        \begin{quote}
        Table 1 reports, for a range of districts, the dimensions of the matrix $\boldsymbol{B}$, the number of constraints returned, the number of nontrivial constraints, and the runtime. The cost is dominated by the V-to-H conversion. Before duplicates are removed, $\boldsymbol{B}$ has one column for each joint value of the response functions, so $\prod_{V \in \bV_1} |V|^{k_V}$ columns, where $k_V$ is the number of configurations of the observed parents of $V$. This number grows exponentially in the number of parent configurations, and so doubly exponentially in the number of parents, and the number of facets can grow faster still. Different joint values of the response functions can give the same column, and these are removed before the V-to-H conversion; this reduces the ternary IV problem from 729 to 333 columns. By contrast, identifying the nontrivial constraints with Proposition 6 takes negligible time.
        \end{quote}
    \end{mdframed}



    \item The screening procedure on pp.~16--17 can avoid listing every vertex. \ldots\ $\max_{p \in \prod_b \Delta_b} h^\top p = \sum_b \max_j h_{bj}$.
    \begin{mdframed}
        We thank the reviewer for this observation, which is correct: a row $h^\top \boldsymbol{p} \le c$ holds for every $\boldsymbol{p} \in \boldsymbol{\check{P}}$ if and only if $\sum_b \max_j h_{bj} \le c$. While implementing it, we found that this check over the product of simplices gives an answer that depends on how the row happens to be written, so it can only flag constraints as potentially nontrivial. Section 3.3 explains this with an example:
        \begin{quote}
        Part (ii) must take the equality constraints into account. A row of $H_i$ is determined only up to a positive multiple and the addition of a linear combination of the rows of $H_e$, and which representative the double description method returns depends on its input. For example, in the graph $X \rightarrow A \leftarrow U \rightarrow B \leftarrow Y$ with binary variables, writing $p(ab \mid xy)$ for $\Pdo(A = a, B = b \mid X = x, Y = y)$, one of the returned rows is
        \[
        p(00 \mid 11) \leq p(00 \mid 10) + p(01 \mid 10).
        \]
        It is violated at points of $\boldsymbol{\check{P}}$ with $p(00 \mid 11) = 1$ and $p(00 \mid 10) = p(01 \mid 10) = 0$. But adding the equality constraint $p(00 \mid 10) + p(01 \mid 10) = p(00 \mid 11) + p(01 \mid 11)$, which is of type (O) with $\bS = \{A\}$, turns it into $p(01 \mid 11) \geq 0$.
        \end{quote}
        In this district, adding random combinations of the equality rows to the inequality rows changes the number of rows flagged by the closed form from 15 to 24, while the number identified by Proposition 6 stays at 8 (script \texttt{14-nontrivial-criterion.R}). We therefore went further and replaced the screening rule with an exact criterion. Definition 4 defines trivial constraints, and Proposition 6 characterizes them:
        \begin{quote}
        \begin{enumerate}
        \item[(ii)] An inequality constraint $h^\top \boldsymbol{p} \leq c$ in $H_i$ is trivial if and only if some entry $p_k$ of $\boldsymbol{p}$ is zero at every column $\beta$ of $\boldsymbol{B}$ with $h^\top \beta = c$. In that case $c - h^\top \boldsymbol{p} = \alpha p_k + \lambda^\top (H_e \boldsymbol{p} - b_e)$ for some $\alpha > 0$ and vector $\lambda$, so the constraint is the nonnegativity constraint $p_k \geq 0$ written in a different form. If the constraint is nontrivial, then there is a strictly positive distribution of the observed variables that satisfies the conditional independencies and every constraint of every district except this one, which it violates.
        \end{enumerate}
        \end{quote}
        Part (i) gives the corresponding result for equality constraints. In keeping with the reviewer's point, the criterion needs no vertex enumeration or optimization:
        \begin{quote}
        Both parts of Proposition 6 are exact and require no optimization. Part (i) needs only a few rank computations. Part (ii) needs only the incidence between facets and columns of $\boldsymbol{B}$: if $I$ is the 0/1 matrix with $I_{r\beta} = 1$ when column $\beta$ of $\boldsymbol{B}$ lies on the facet defined by row $r$ of $H_i$, then row $r$ is nontrivial if and only if every entry of row $r$ of $I \boldsymbol{B}^\top$ is positive.
        \end{quote}
        The software now implements this criterion. The tripartite Bell district takes 26 seconds (Table 1), where the listing of extreme points in the submitted version would have required $8^8 \approx 1.7\times 10^7$ of them.
    \end{mdframed}

    \item The displayed equality on p.~24 cannot hold as written. \ldots\, Please correct the expression and state the four intended equality constraints explicitly.
    \begin{mdframed}

Thank you; this was a typo and we regret the error. The constraints are now stated explicitly (Section 5):
\begin{quote}
Besides normalization, the equality constraints are
\[
\Pdo(V_1 = v_1, V_6 = v_6 \mid V_2 = 0) = \Pdo(V_1 = v_1, V_6 = v_6 \mid V_2 = 1), \quad v_1, v_6 \in \{0, 1\},
\]
which are of type (O) with $\bS = \{V_1, V_6\}$.
\end{quote}
Given normalization, three of these four are linearly independent, which is why the district has 5 linearly independent equality constraints in total.
    \end{mdframed}
\end{enumerate}


\section*{Reviewer 2}

\begin{enumerate}

    \item[1.] \ldots\, is there a way to relate nontrivial inequalities to specific graphical structures, such as exclusion restriction?

    \item[3.] \ldots\, can such differences be characterized by some graph criteria or rules? \ldots


    \begin{mdframed}
    We respond to these two related comments together. Both are good questions, and the revision adds a new subsection (Section 4.3) with a new remark (Remark 3), and a paragraph in the Discussion.


\emph{A graphical criterion.} Remark 3 gives two structures that guarantee nontrivial inequality constraints:
\begin{quote}
Let $\mathcal{G}$ satisfy Conditions 1 and 2 and have c-degree one, and let $D$ be a district of $\mathcal{G}$ whose observed members and external parents each have at least two levels, and that contains at least one of the following:
\begin{enumerate}
\item[(i)] (\emph{instrumental type}) members $A$ and $B$, with $A$ a parent of $B$, and a directed path from some $W \in \bV_2$ to $A$, all of whose vertices other than $W$ are members of $D$, such that no vertex on the path other than $A$ is a parent of $B$;
\item[(ii)] (\emph{Bell type}) members $A$ and $B$ and external parents $W_A, W_B \in \bV_2$ such that $W_A$ reaches $A$ but not $B$, and $W_B$ reaches $B$ but not $A$.
\end{enumerate}
Then the H-representation of $D$ contains a nontrivial inequality constraint.
\end{quote}
This relates the inequalities to exclusion restrictions, as the reviewer suggests, and the manuscript places each of our examples:
\begin{quote}
In both types, the inequality constraints come from exclusion restrictions. In type (i), $W$ and the members on the path to $A$ are excluded from $B$, so $W$ acts as an instrument for the effect of $A$ on $B$. In type (ii), $W_A$ is excluded from $B$ and $W_B$ from $A$. Type (i) appears in the instrumental variable model, in both districts of Figure 3, and in the right graph of Figure 5, with $(W, A, B) = (V_1, V_2, V_4)$. Type (ii) appears in the tripartite Bell district, in $X \rightarrow A \leftarrow U \rightarrow B \leftarrow Y$, and in the district $\{B, D\}$ of $A \rightarrow B \rightarrow C \rightarrow D$ with $U \rightarrow B$ and $U \rightarrow D$. The merged second district of Figure 7 contains both. The districts in our examples without nontrivial inequality constraints contain neither: the front door model, the left graph of Figure 5, and the first district of Figure 7.
\end{quote}

\emph{Comparing candidate graphs.} This addresses comment 3 as follows (Section 7):
\begin{quote}
Because the constraints are derived one district at a time, two candidate graphs need only be compared on the districts where they differ, either in membership or in the parent sets of members. All other districts and their constraints are shared. The comparison also does not require deriving the full constraint set for every candidate graph. To find out whether a constraint of graph $\mathcal{G}_2$ distinguishes it from $\mathcal{G}_1$, one can sample distributions from $\mathcal{M}(\mathcal{G}_1)$ and check whether they violate it, as illustrated for Figure 5 in the reproducibility scripts in the Supplementary Material. Remark 3 also shows where to look. If a district of one candidate graph contains an instrumental-type or Bell-type structure, it has nontrivial inequality constraints, and these can distinguish it from a candidate graph whose corresponding district contains neither structure. In Figure 5, the district $\{V_2, V_4\}$ is of instrumental type in the right graph and contains neither structure in the left graph (Section 4.3), and the instrumental inequalities of the right graph distinguish the two.
\end{quote}
The script \texttt{06-nested-equivalent.R} confirms this for Figure 5: distributions sampled from the left graph violate exactly the 8 nontrivial constraints of the right graph, and never violate those of the left graph.

\emph{Limitations and related work.} Remark 3 is a sufficient condition only. When a district contains neither structure, it does not show that the district has no inequality constraints, and the Algorithm or the sampling check above is still needed, unless the criterion of \citet{henson2014theory} applies. The manuscript sets this in the context of existing work on when a graph implies inequality constraints:
\begin{quote}
Remark 3 gives a sufficient condition only, and these mechanisms are known. Whether a graph implies any inequality constraints at all has been studied for hidden variable DAGs in general. \citet{evans2023latent} shows that a discrete hidden variable DAG implies no nontrivial inequality constraints if and only if its marginal model is that of a DAG without hidden variables, and gives the CHSH inequalities of the district $\{B, D\}$ above as an example. \citet{khanna2024classifying} use this result, together with d-separation, e-separation and incompatible supports, to classify graphs with up to four observed variables, leaving 3 of the 2809 graphs with four observed variables unresolved. Graphs with a structure of type (i) are covered by their result on nonmaximal graphs: $W$ and $B$ are not adjacent, and they cannot be d-separated, because conditioning on a member on the path from $W$ to $B$ opens a path through the shared unobserved variable. In the other direction, if edges can be added using Proposition 1 until the unobserved variables can be removed without changing the conditional independencies, then the graph implies no inequality constraints \citep{henson2014theory}. Whether this criterion is also necessary is the conjecture discussed by \citet{khanna2024classifying}, which is open in general.
\end{quote}
Remark 3 also does not say which inequalities arise:
\begin{quote}
It identifies structures that produce inequality constraints, but not which inequalities arise. These form many families; for example, instrumental-type districts acquire new inequalities when the instrument has three or more levels \citep{bonet2001instrumentality}.
\end{quote}
So the exact form of the distinguishing inequalities still requires our algorithm. We also relate the result to e-separation:
\begin{quote}
Remark 3 is also related to the e-separation criterion of \citet{evans2012graphical}. Sets $A$ and $B$ are e-separated given $C$ after deleting $D$ if they are d-separated given $C$ in the graph with the nodes in $D$ removed. If, in addition, no vertex in $C$ is a descendant of a vertex in $D$, then for each value $d$ of $D$, the probabilities $\Pr(a, b, d \mid c)$ must be compatible with a distribution in which $A \indep B \mid C$ \citep[Theorem 4.2]{evans2012graphical}. This is often an inequality constraint \citep{evans2012graphical, finkelstein2021entropic}, but it need not be a new constraint: in the left graph of Figure 5, $V_1$ and $V_3$ are e-separated after deleting $V_2$, which only restates $V_1 \indep V_3 \mid V_2$. e-separation applies to graphs of any c-degree, but it gives only a sufficient condition, and it detects inequality constraints of the instrumental type: in the instrumental variable model, $Z$ and $Y$ are e-separated after deleting $X$, which yields the instrumental inequality, and in the right graph of Figure 5, $V_1$ and $V_4$ are e-separated after deleting $\{V_2, V_3\}$. It does not detect those of the Bell type: \citet[Remark 2.6]{evans2023latent} notes that the CHSH inequality cannot be read off in this way. In the graph $X \rightarrow A \leftarrow U \rightarrow B \leftarrow Y$, for example, $A$ and $B$ share the unobserved $U$, which cannot be deleted, while the pairs that can be separated, such as $X$ and $B$, are marginally independent, which yields only that independence. Remark 3 is restricted to districts of c-degree one but covers both types.
\end{quote}
    \end{mdframed}

   \item[2.] \ldots\, is there a method to remove redundant constraints that can be derived from other constraints \ldots?
    \begin{mdframed}
        This is addressed in the Discussion (Section 7):
        \begin{quote}
        Within a district, the H-representation returned by the double description method is irredundant: each inequality defines a facet of the district polytope, and the equality constraints form a basis \citep{fukuda1996double}. By Proposition 6(ii), a nontrivial inequality is also not implied by the conditional independencies together with all other constraints, including those of other districts. So the only constraints that can be removed are the trivial ones, which Proposition 6 identifies. For large problems, the number of constraints can be reduced further by grouping them into orbits under the symmetries of the model, such as relabelling values of a variable or permuting exchangeable nodes. For example, the 53,856 tripartite Bell inequalities fall into 46 symmetry classes, one of which is nonnegativity \citep{Sliwa2003Symmetries}.
        \end{quote}
        As a check, applying linear-programming redundancy removal (\texttt{rcdd::redundant}) to every example in Table 1 except the tripartite Bell district removed no constraints.

        As we were investigating this, we did discover that the V-representations that we construct may sometimes contain redundant columns. We implemented and tested a simple modification to remove those, and it sped up the subsequent V-to-H conversion, which yielded identical results in all of our examples and test cases. This is described in Section 6:
        \begin{quote}
        Different joint values of the response functions can give the same column, and these are removed before the V-to-H conversion; this reduces the ternary IV problem from 729 to 333 columns.
        \end{quote}
    \end{mdframed}


    \item[4.] \ldots\, it would be better to add a brief implementation section \ldots\, demonstrate the corresponding R commands to obtain a complete set of constraints.
    \begin{mdframed}
        We have added Section 6, ``Implementation''. It shows the R commands for the IV model:
\begin{verbatim}
library(meraconstraints)
ivdag  <- caugi(Z %-->% X, X %-->% Y, U %-->% X + Y) |> dhvm()
ivmera <- meraconstraints(ivdag)
ivmera                        # summary of the constraints
H <- ivmera$polytope$Hrep[[1]]
P <- ivmera$functional_relations$probform[[1]]
## all constraints, written in observed probabilities
hrep_to_text(H, P)
## only the nontrivial ones
hrep_to_text(H, P)[ivmera$nontrivial[[1]]]
## test a given vector of district factors
hrep_test(H, c(0.101, 0.338, 0.013, 0.548, 0.317, 0.51, 0.004, 0.169),
          which = TRUE)
\end{verbatim}
        It also covers graphs with c-degree greater than one and the supplementary material:
        \begin{quote}
        For a graph with c-degree greater than one, \texttt{dhvm\_merge\_latents()} carries out the merging step of Section 5 before \texttt{meraconstraints()} is called. R scripts that reproduce every example, number and table in this paper, including the computational benchmarks of Table 1, are provided in the supplementary material.
        \end{quote}
        The scripts come with a README describing which script corresponds to which section.
    \end{mdframed}
    \item[5.] \ldots\, whether the proposed method can be extended to settings where some variables in $\mathbf{V}$ are continuous? \ldots
    \begin{mdframed}
        The Discussion (Section 7) now explains that two kinds of continuous variables can be handled, and it agrees with the reviewer about the third:
        \begin{quote}
        It relies on each member of a district having finitely many response functions, but two kinds of continuous variables can be handled. First, if a continuous variable $V_c$ has no observed children, as for a continuous outcome, and $g$ maps its range onto the cells of a finite partition, then $g(V_c)$ is a function of the same parents and no other structural equation changes. So $g(V_c)$ satisfies every constraint our method derives for the graph with a discrete variable in place of $V_c$. The partition-based inequalities of \citet{kedagni2020generalized} for a continuous outcome are of this type. Second, if the members of a district are discrete but an external parent $Z$ is continuous, as for a continuous instrument, we can apply the method treating $Z$ as if it took $k$ levels $z_1, \ldots, z_k$. The derivation of each resulting constraint does not depend on which values these labels stand for, so it holds for every choice of $k$ values of $Z$, or equivalently for the supremum of its left-hand side over such choices. The inequalities of \citet{kedagni2020generalized} for a continuous instrument have this form. Partitioning the range of an exogenous $Z$ and treating the cell as a categorical instrument also gives valid constraints, but weaker ones, because averaging within a cell can hide a violation between two values of $Z$ in the same cell. If a continuous variable is a parent of another member of its own district, the space of response functions is infinite-dimensional and the polytope argument does not apply. Discretizing it does not give valid constraints in general: the instrumental variable model with a continuous treatment has no testable implications \citep{gunsilius2021nontestability}, whereas the model with the treatment discretized implies the instrumental inequalities.
        \end{quote}
        We leave this last case to future work.
    \end{mdframed}
\end{enumerate}


\bibliographystyle{abbrvnat}
\bibliography{references.bib}


\end{document}
